\section{Entropy and identification of the limit}\label{sec:entropy}

The ball estimate has exponent $n$, whereas $A_n$ has dimension $n-1$.
We follow the ordinary-ball strategy of
\cite[Section~2.7.1]{ELMV2011}. We retain the estimate for every invariant
probability dominated by a finite multiple of the limit, so it forces
positive entropy after restriction to any positive-weight collection of
ergodic components. Measure classification then identifies the limit.

For an invariant probability $\nu$, write $h_\nu(a)$ for the
measure-theoretic entropy of the map $x\mapsto xa$.
Fix distinct real numbers $\lambda_1,\ldots,\lambda_n$ with sum zero and put
\[
 a(t)=\operatorname{diag}(e^{\lambda_1t},\ldots,e^{\lambda_nt}),
 \qquad \lambda=\min_{i\ne j}|\lambda_i-\lambda_j|>0.
\]

\begin{proposition}[Entropy from ordinary balls]\label{prop:entropy-from-balls}
Let $n\ge3$ and let $\mu$ be an $A_n$-invariant probability on $X_n$.
Suppose that for every compact $\Omega\subset X_n$ there are
$C_\Omega,r_\Omega>0$ such that
\[
 \mu(xB(r))\le C_\Omega r^n
 \qquad(x\in\Omega,\ 0<r<r_\Omega).
\]
Then every $A_n$-invariant probability $\xi$ satisfying $\xi\le c\mu$
for some finite $c$ has $h_\xi(a(1))\ge\lambda/3$.
\end{proposition}

\begin{proof}
Fix $0<b<1/4$ and consider the two-sided neighborhoods
\[
 T_t=a(t)B(b)a(-t)\cap a(-t)B(b)a(t),\qquad t\ge0.
\]
If $h\in T_t$, then
\begin{equation}\label{eq:tube-entries}
 |h_{ii}-1|<b,\qquad
 |h_{ij}|<b e^{-|\lambda_i-\lambda_j|t}\quad(i\ne j).
\end{equation}
Set $r=e^{-\lambda t}$.  Every term in the determinant expansion other
than the product of the diagonal entries contains at least two
off-diagonal entries.  Since $\det h=1$, it follows that
$\prod_i h_{ii}=1+O_{n,b}(r^2)$.  Thus
\[
 d(h)=\operatorname{diag}\left(h_{11},\ldots,h_{n-1,n-1},
                   \Bigl(\prod_{i<n}h_{ii}\Bigr)^{-1}\right)
\]
belongs to a fixed compact subset $\mathcal D\subset A_n$, and
$d(h)^{-1}h\in B(Cr)$ for a constant depending only on $n,b$.

In logarithmic coordinates, $\mathcal D$ is a bounded subset of $\R^{n-1}$.
Covering it by a mesh of size comparable to $r$, and enlarging $C$ to
allow for matrix multiplication, gives
\begin{equation}\label{eq:tube-cover}
 T_t\subset\bigcup_{j=1}^{N_t}d_jB(Cr),
 \qquad d_j\in\mathcal D,\qquad N_t\le Cr^{-(n-1)}.
\end{equation}
For $x$ in a compact set $\Omega$, all the centers $xd_j$ lie in the
fixed compact set $\Omega\mathcal D$.  The assumed ball estimate and domination
therefore imply, for all sufficiently large $t$,
\begin{equation}\label{eq:tube-mass}
 \xi(xT_t)\le C_{\Omega,b,c}\,r^{-(n-1)}r^n
             =C_{\Omega,b,c}e^{-\lambda t}
 \qquad(x\in\Omega).
\end{equation}
Since $\xi$ is a probability, the constant can be enlarged to cover
the finitely many integer times before this estimate holds.

We use the entropy criterion of
Einsiedler--Lindenstrauss--Michel--Venkatesh
\cite[Corollary~3.3]{ELMV2009}.  In the notation
\[
 B^{(s,t)}=a(-s)Ba(s)\cap a(-t)Ba(t),
\]
the criterion states the following: if flow-invariant probabilities
$\xi_i$ converge weakly to a probability $\xi$, and $t_i\to\infty$,
then $h_\xi(a(1))\ge\eta$ provided that, for every compact set $\Omega$,
there are an identity neighborhood $B$ and a constant $C_\Omega$ such that
\[
 \xi_i\bigl(xB^{(-t_i,t_i)}\bigr)
       \le C_\Omega e^{-2\eta t_i}\qquad(x\in\Omega).
\]
This criterion applies on the noncompact lattice space; its requirement
that the limit be a probability is essential.  Here take the constant
sequence $\xi_i=\xi$, the neighborhood $B(b)$, $t_i=i$, and
$\eta=\lambda/3$.  Since $B(b)^{(-t,t)}=T_t$,
Equation~\eqref{eq:tube-mass} proves the required estimate.
\end{proof}

\begin{lemma}[Entropy in the ergodic components]\label{lem:entropy-components}
Under the hypotheses of Proposition~\ref{prop:entropy-from-balls},
almost every $A_n$-ergodic component of $\mu$ has positive entropy
for $a(1)$.
\end{lemma}

\begin{proof}
The group $A_n\simeq\R^{n-1}$ is locally compact and second countable,
and its action on the standard Borel space $X_n$ is continuous.
The ergodic decomposition theorem
\cite[Theorem~5.2]{GreschonigSchmidt2000}, specialized to invariant
probabilities, therefore supplies a measurable decomposition
\[
 \mu=\int_Z\nu_z\,\dd\tau(z)
\]
into $A_n$-invariant, $A_n$-ergodic probabilities over a factor fixed
by $A_n$. Set $T(x)=xa(1)$.
The function $z\mapsto h_{\nu_z}(T)$ is measurable.  Indeed, choose
increasing finite Borel partitions $\mathcal P_k$ generating the Borel
$\sigma$-algebra of $X_n$.  For every $T$-invariant probability $\nu$,
\[
 h_\nu(T)=\sup_k\inf_{N\ge1}\frac1N
 H_\nu\left(\bigvee_{j=0}^{N-1}T^{-j}\mathcal P_k\right),
\]
where $H_\nu$ denotes the Shannon entropy of a finite partition.
The right side is a countable combination of Borel functions of $\nu$.

The component partition is fixed by $T$, since $a(1)\in A_n$.
The entropy integral formula for a fixed invariant partition
\cite[Theorem~9.8]{Rokhlin1967} consequently gives
\begin{equation}\label{eq:entropy-decomposition}
 h_\mu(T)=\int_Z h_{\nu_z}(T)\,\dd\tau(z).
\end{equation}
The same formula holds after restricting to any measurable collection
of components.  It does not require the $\nu_z$ to be $T$-ergodic:
the induced transformation on the component space is the identity
and contributes zero entropy.

If $Z_0=\{z:h_{\nu_z}(T)=0\}$ had weight $w=\tau(Z_0)>0$, then
\[
 \xi=\frac1w\int_{Z_0}\nu_z\,\dd\tau(z)
\]
would be an $A_n$-invariant probability with $\xi\le w^{-1}\mu$.
Equation~\eqref{eq:entropy-decomposition}, applied to this restricted
decomposition, would give $h_\xi(T)=0$, contradicting
Proposition~\ref{prop:entropy-from-balls}.
\end{proof}

\paragraph{Measure classification.}
Einsiedler--Katok--Lindenstrauss
\cite[Theorem~1.3 and Corollary~1.4]{EKL2006} prove that, for prime
$n\ge3$, an $A_n$-invariant and $A_n$-ergodic probability on
$\operatorname{SL}_n(\R)/\operatorname{SL}_n(\Z)$ is Haar whenever
some diagonal one-parameter subgroup has positive entropy.
Their action is on the left.  The map
$\Gamma g\mapsto g^{-1}\Gamma$, where $\Gamma=\operatorname{SL}_n(\Z)$,
changes our right action by $a(t)$ into the left action by $a(-t)$.
Since entropy is unchanged by inversion of an invertible transformation,
their result applies to our convention without any additional hypothesis.

\begin{proof}[Proof of Theorem~\ref{thm:main}]
Let $\mu_i$ be the packet probabilities for a sequence in the theorem.
Equation~\eqref{eq:discriminant-scale} gives $Q_i\to\infty$.
By Proposition~\ref{prop:tightness}, this sequence is tight and every
weak subsequential limit $\mu$ is an $A_n$-invariant probability.
Proposition~\ref{prop:ball-bound} supplies the hypotheses of
Proposition~\ref{prop:entropy-from-balls}.
Lemma~\ref{lem:entropy-components} shows that almost every
$A_n$-ergodic component of $\mu$ has positive entropy for $a(1)$.
As $n$ is prime, the classification just stated makes every such
component equal to $m_n$.  Integrating the components gives $\mu=m_n$.

Thus every convergent subsequence has the same probability limit.
Together with tightness, this proves $\mu_i\to m_n$ and excludes
loss of mass.  All preceding estimates were for each prescribed local
homothety packet separately, so no averaging over different local
module types is used.
\end{proof}
